% Pregunta 949 (2002-octubre-OEP2002, n.º 32): Dornbusch con expansión monetaria anunciada en t0
% y aplicada en t1. Arriba, el tipo de cambio: salta en t0, sigue depreciándose hasta t1 (c),
% sobrerreacciona por encima de su nuevo nivel de largo plazo y luego se aprecia (d es falsa).
% Abajo, el tipo de interés: sube por encima de i* entre t0 y t1 y cae por debajo en t1 (b es falsa).
% Autor: Víctor Gutiérrez Marcos
\documentclass[tikz,border=6pt]{standalone}
\input{tikz-tcee}
\begin{document}
\begin{tikzpicture}[tcee]
  % ---------- Panel superior: tipo de cambio ----------
  \begin{scope}
    \fill[area] (1.5,0) rectangle (3.6,4.3);
    \ejes{$t$}{$e$}{6.4}{4.6}
    \draw[guia] (1.5,0) -- (1.5,4.3);
    \draw[guia] (3.6,0) -- (3.6,4.3);
    \node[etiqueta,below] at (1.5,0) {$t_0$};
    \node[etiqueta,below] at (3.6,0) {$t_1$};
    % Niveles de largo plazo
    \draw[guia] (0,2.7) -- (6.2,2.7);
    \node[etiqueta,left] at (0,0.9) {$\bar e_0$};
    \node[etiqueta,left,text=tceerojo] at (0,2.7) {$\bar e_1$};
    % Senda
    \draw[curva1] (0,0.9) -- (1.5,0.9);
    \draw[curvanueva,densely dotted] (1.5,0.9) -- (1.5,1.8);
    \draw[curvanueva] (1.5,1.8) .. controls (2.6,1.8) and (3.2,2.6) .. (3.6,3.4)
      .. controls (4.2,2.9) and (5,2.72) .. (6.2,2.7);
    \node[punto] at (3.6,3.4) {};
    % Anotaciones
    \node[etiqueta,font=\footnotesize,align=left,anchor=west] at (1.6,1.15) {salto\\en $t_0$};
    \node[etiqueta,font=\footnotesize,text=tceeverde,align=center] at (2.55,3.85) {c) se deprecia\\en $t_0$ y hasta $t_1$};
    \node[etiqueta,font=\footnotesize,align=left,anchor=south west] at (3.7,3.4) {sobrerreacción};
    \node[etiqueta,font=\footnotesize,text=tceerojo,align=left,anchor=north west] at (4.0,2.55) {desde $t_1$ se aprecia\\(d es falsa)};
  \end{scope}
  % ---------- Panel inferior: tipo de interés ----------
  \begin{scope}[yshift=-4.4cm]
    \fill[area] (1.5,0) rectangle (3.6,3.1);
    \ejes{$t$}{$i$}{6.4}{3.4}
    \draw[guia] (1.5,0) -- (1.5,3.1);
    \draw[guia] (3.6,0) -- (3.6,3.1);
    \node[etiqueta,below] at (1.5,0) {$t_0$};
    \node[etiqueta,below] at (3.6,0) {$t_1$};
    \draw[guia] (0,1.3) -- (6.2,1.3);
    \node[etiqueta,left] at (0,1.3) {$i^*$};
    % Senda
    \draw[curva1] (0,1.3) -- (1.5,1.3);
    \draw[curvanueva] (1.5,1.3) .. controls (2.5,1.35) and (3.1,1.8) .. (3.6,2.3);
    \draw[curvanueva,densely dotted] (3.6,2.3) -- (3.6,0.45);
    \draw[curvanueva] (3.6,0.45) .. controls (4.3,1.0) and (5,1.28) .. (6.2,1.3);
    % Anotaciones
    \node[etiqueta,font=\footnotesize,align=center] at (2.5,2.55) {$i>i^*$:\\se espera\\depreciación};
    \node[etiqueta,font=\footnotesize,text=tceerojo,align=left,anchor=west] at (3.75,2.35) {en $t_1$, $i$ baja\\(b es falsa)};
    \node[etiqueta,font=\footnotesize,align=left,anchor=north west] at (4.3,0.95) {$i<i^*$: se espera\\apreciación};
  \end{scope}
\end{tikzpicture}
\end{document}
